\section{Signed-Schur verification on two complete graph branches}
\label{sec:signed-schur-details}

This appendix supplies the derivative conditions and proof for the
verification procedure in Section~\ref{sec:schur-seam-verifier}.  Its input is
a pair of graph parametrizations.  Tangential elimination reduces coverage
and uniqueness to a scalar growth test with opposite signs on the two
branches.  The resulting estimates provide the J-side inequality; the
error bound is verified separately.

Fix \(\lambda>0\).  For graph points \(g=(u,u^*)\), use the algebraic
coordinates
\[
 \mathsf Xg=u+\lambda u^*,\qquad
 \mathsf Yg=u-\lambda u^*.
\]
Choose orthogonal input coordinates
\(\mathbb R^n=\mathbb R^{n-1}\times\mathbb R\), with projections
\(P\) and \(N\).  Norms below are Euclidean operator norms.

\begin{theorem}[Signed Schur-growth certificate]
\label{thm:signed-schur-growth}
Let \(n\ge2\) and \(R,T,\mu>0\).
Fix nonnegative constants \(H,B,K_z,K_+,K_-\).
Suppose that two \(C^1\) graph parametrizations
\[
 g_\sigma(a,t)=(u_\sigma(a,t),v_\sigma(a,t)),
 \qquad (a,t)\in\overline B_R^{n-1}\times[0,T],
 \quad \sigma\in\{+1,-1\},
\]
extend as \(C^1\) maps to a neighborhood of their displayed parameter
domains, and their restrictions to those displayed domains take values in
\(\operatorname{gph}F\).  No graph membership is required for the auxiliary
extension outside those domains.  They agree on the seam:
\[
 g_+(a,0)=g_-(a,0)=g_0(a).
 \tag{S1}\label{eq:schur-seam}
\]
Write \(X_\sigma=\mathsf Xg_\sigma\),
\(Y_\sigma=\mathsf Yg_\sigma\), and assume
\(PX_\sigma(0,0)=0\).  At every displayed parameter point, form
\[
 A_\sigma=P D_aX_\sigma,
 \qquad b_\sigma=P\partial_tX_\sigma.
\]
Assume the following bounds hold uniformly on both complete parameter
domains:
\begin{align}
 \operatorname{sym} A_\sigma&\succeq\mu I,
 &\|b_\sigma\|&\le H,
 \tag{S2}\label{eq:schur-tangent}\\
 \|N D_aX_\sigma A_\sigma^{-1}\|&\le B,
 \tag{S3}\label{eq:schur-normal-tangent}\\
 \|D_aY_\sigma A_\sigma^{-1}\|&\le K_z,
 \tag{S4}\label{eq:schur-output-tangent}\\
 \|\partial_tY_\sigma-
       D_aY_\sigma A_\sigma^{-1}b_\sigma\|&\le K_\sigma.
 \tag{S5}\label{eq:schur-output-normal}
\end{align}
Let \(h_+,h_-:[0,T]\to[0,\infty)\) be \(C^1\), convex, strictly
increasing, and zero at zero.  The signed normal Schur derivatives are
required to satisfy
\[
 \sigma\bigl(N\partial_tX_\sigma-
       ND_aX_\sigma A_\sigma^{-1}b_\sigma\bigr)
       \ge h_\sigma'(t).
 \tag{S6}\label{eq:schur-normal-growth}
\]
Choose \(0<r<\mu R-HT\).  Then the following conclusions hold.
\begin{enumerate}
\item For every \(z\in B_r^{n-1}\) and \(t\in[0,T]\), there is exactly
one \(a=\alpha_\sigma(z,t)\in B_R^{n-1}\) with
\(PX_\sigma(a,t)=z\).  The seam function
\[
 \nu(z)=NX_\sigma(\alpha_\sigma(z,0),0)
\]
is independent of \(\sigma\).  The represented graph covers the open
input neighborhood
\[
 U=\{(z,s):z\in B_r^{n-1},\quad
       \nu(z)-h_-(T)<s<\nu(z)+h_+(T)\}.
 \tag{S7}\label{eq:schur-covered-neighborhood}
\]
\item On the represented graph with tangential inputs in \(B_r^{n-1}\),
\(\mathsf X\) has no collisions: equal inputs imply equal graph points.
At the seam, the two parameter representations describe the same graph
point, not two different solutions.
\item Set \(C_B=\sqrt{1+B^2}\), and define
\[
 V(v)=\max\{K_+s+K_-t:0\le s,t\le T,
                         \ h_+(s)+h_-(t)\le v\},\qquad v\ge0.
 \tag{S8}\label{eq:schur-allocation}
\]
For every pair of represented graph points \(g=(u,v)\) and
\(g'=(u',v')\) whose tangential inputs belong to \(B_r^{n-1}\),
\[
 \begin{gathered}
 \|u-u'\|^2+\lambda^2\|v-v'\|^2
 \le\frac12\bigl[d^2+\{K_zd+V(C_Bd)\}^2\bigr],\\
 d=\|\mathsf Xg-\mathsf Xg'\|,
 \end{gathered}
 \tag{S9}\label{eq:schur-RL-gauge}
\]
holds.  The bound \(\omega(d)=K_zd+V(C_Bd)\) is nondecreasing and
continuous at zero.  In the power case below it gives explicit RL
parameters.  On the covered input neighborhood, the left-hand side is
\(\|Jx-Jx'\|^2+\|(x-Jx)-(x'-Jx')\|^2\).
\end{enumerate}
If the intended graph localization \(\mathcal G\) satisfies the complete
fiber identity
\[
 \mathcal G\cap\mathsf X^{-1}(U)
 =\left(\bigcup_{\sigma=\pm1}
    g_\sigma(\overline B_R^{n-1}\times[0,T])\right)
      \cap\mathsf X^{-1}(U),
 \tag{S10}\label{eq:schur-complete-fibers}
\]
then the conclusions over covered inputs apply to
\(\mathcal G\cap\mathsf X^{-1}(U)\), including a
single-valued localized resolvent on all of \(U\).  For an unrestricted
resolvent conclusion, \(\mathcal G\) must be the complete graph of \(F\)
over these inputs.
\end{theorem}

\begin{proof}
Fix \(t\) and \(\sigma\), and put \(f(a)=PX_\sigma(a,t)\).
Integration of \eqref{eq:schur-tangent} along a line segment gives
\[
 \langle f(a)-f(a'),a-a'\rangle\ge\mu\|a-a'\|^2.
 \tag{S11}\label{eq:schur-tangent-monotone}
\]
Also \(\|f(0)\|\le HT\).  For \(\|a\|=R\), \(\|z\|<r\),
\[
 \langle f(a)-z,a\rangle
 \ge\mu R^2-(HT+r)R>0.
\]
The continuous map \(f-z\) therefore has a zero in the ball.  Indeed, if
it did not, the self-map
\(a\mapsto-R(f(a)-z)/\|f(a)-z\|\) of the closed ball would have a
Brouwer fixed point on the sphere, contradicting the displayed positive
inner product.  Uniqueness follows from
\eqref{eq:schur-tangent-monotone}.  Since \(A_\sigma\) is invertible,
the parameterized implicit function theorem yields a \(C^1\) map
\(\alpha_\sigma\), with
\[
 D_z\alpha_\sigma=A_\sigma^{-1},\qquad
 \partial_t\alpha_\sigma=-A_\sigma^{-1}b_\sigma.
 \tag{S12}\label{eq:schur-partial-inverse}
\]
All matrices on the right are evaluated at the corrected parameter
\((\alpha_\sigma(z,t),t)\).  At the seam,
\[
 \alpha_+(z,0)=\alpha_-(z,0).
\]

Define, solely for this proof,
\[
 n_\sigma(z,t)=NX_\sigma(\alpha_\sigma(z,t),t),\qquad
 y_\sigma(z,t)=Y_\sigma(\alpha_\sigma(z,t),t).
\]
The chain rule, \eqref{eq:schur-partial-inverse}, and the hypotheses give
\[
 \|D_zn_\sigma\|\le B,\quad
 \sigma\partial_t n_\sigma\ge h_\sigma'(t),\quad
 \|D_zy_\sigma\|\le K_z,\quad
 \|\partial_ty_\sigma\|\le K_\sigma.
 \tag{S13}\label{eq:schur-reduced-bounds}
\]
Both normal coordinates start at \(\nu(z)\).  Integration shows that
\(n_+(z,\cdot)\) is strictly increasing and reaches at least
\(\nu(z)+h_+(T)\); \(n_-(z,\cdot)\) is strictly decreasing and reaches
at most \(\nu(z)-h_-(T)\).  The intermediate value theorem proves
\eqref{eq:schur-covered-neighborhood}.  Their ranges meet only at the
seam, and each is one-to-one.  Consequently equal full Minty inputs
imply the same branch parameter, except for the duplicate seam
representation, where \eqref{eq:schur-seam} gives the same graph point.

For convex \(h\) with \(h(0)=0\),
\[
 h(t)-h(s)\ge h(t-s),\qquad 0\le s\le t\le T.
 \tag{S14}\label{eq:schur-convex-increment}
\]
Consider two points on the same branch with reduced parameters
\((z,t)\) and \((w,s)\).  Equations
\eqref{eq:schur-reduced-bounds}--\eqref{eq:schur-convex-increment} give
\[
 h_\sigma(|t-s|)
 \le|n_\sigma(z,t)-n_\sigma(w,s)|+B\|z-w\|
 \le C_Bd.
\]
The difference of the \(Y\)-coordinates is at most
\(K_z\|z-w\|+K_\sigma|t-s|\), which is bounded by
\(K_zd+V(C_Bd)\), since the pair
\((|t-s|,0)\), or \((0,|t-s|)\), is feasible in
\eqref{eq:schur-allocation}.

For points on opposite branches, use the common seam at the same
tangential coordinate to obtain
\[
 h_+(t)+h_-(s)
 \le n_+(z,t)-n_-(z,s)
 \le|n_+(z,t)-n_-(w,s)|+B\|z-w\|
 \le C_Bd.
\]
Compare their \(Y\)-coordinates through \(y_0(z)\) and \(y_0(w)\):
\[
 \|y_+(z,t)-y_-(w,s)\|
 \le K_+t+K_z\|z-w\|+K_-s
 \le K_zd+V(C_Bd).
\]
Thus \(\|\mathsf Yg-\mathsf Yg'\|\le K_zd+V(C_Bd)\) for every
ordered branch pair, including seam endpoints.  The identity
\[
 2\bigl(\|u-u'\|^2+\lambda^2\|v-v'\|^2\bigr)
 =\|\mathsf Xg-\mathsf Xg'\|^2+\|\mathsf Yg-\mathsf Yg'\|^2
\]
gives \eqref{eq:schur-RL-gauge}.
The maximization defining \(V\) is over a nonempty compact set, and
strict increase of both \(h_\sigma\) forces \(V(v)\to0\) as
\(v\downarrow0\).  The complete-fiber conclusion is then immediate.
\end{proof}

\begin{corollary}[Explicit power constants]
\label{cor:schur-power-constants}
In Theorem~\ref{thm:signed-schur-growth}, suppose
\(h_\sigma(t)=c_\sigma t^p\), with \(c_\sigma>0\), \(p>1\).
Put \(p'=p/(p-1)\) and
\[
 C_{p}=\left[
  (K_+c_+^{-1/p})^{p'}+(K_-c_-^{-1/p})^{p'}
 \right]^{1/p'}.
\]
On every represented input block whose tangential inputs lie in
\(B_r^{n-1}\) and whose diameter is at most \(D>0\),
\[
 \boxed{\quad\gamma=1/p,\qquad
 L=K_zD^{1-1/p}+(1+B^2)^{1/(2p)}C_p.\quad}
 \tag{S15}\label{eq:schur-explicit-L}
\]
For \(p=1\), replace \(C_p\) by
\(\max\{K_+/c_+,K_-/c_-\}\) and take \(\gamma=1\).
If the branch powers differ, \(h_\sigma(t)=c_\sigma t^{p_\sigma}\),
one may take \(p=\max\{p_+,p_-\}\) and replace
\(c_\sigma\) by \(c_\sigma T^{p_\sigma-p}\) in the same formula.
\end{corollary}

\begin{proof}
H\"older's inequality in \(\mathbb R^2\) gives
\(V(v)\le C_pv^{1/p}\), without requiring that the maximizing
parameters remain below \(T\).  Insert this in
\eqref{eq:schur-RL-gauge} and use
\(d\le D^{1-1/p}d^{1/p}\).  The \(p=1\) formula is the elementary
linear-programming bound.  Finally,
\(t^{p_\sigma}\ge T^{p_\sigma-p}t^p\) for \(0\le t\le T\).
\end{proof}

\begin{proposition}[A matched-jet estimate across the seam]
\label{prop:schur-matched-jet}
Suppose the hypotheses of Theorem~\ref{thm:signed-schur-growth} hold
with \(h_+(t)=h_-(t)=ct^p\), \(c>0\), \(p\ge1\).
In the reduced coordinates constructed in its proof, assume the
additional paired-branch estimate
\[
 \|y_+(z,t)-y_-(z,t)\|\le Et^p
 \quad(z\in B_r^{n-1},\ 0\le t\le T).
 \tag{S16}\label{eq:schur-matched-output-jet}
\]
Let \(K=\max\{K_+,K_-\}\).  In the J-side estimate
\eqref{eq:schur-RL-gauge}, the term \(K_zd+V(C_Bd)\) can be replaced by
\[
 \omega(d)=\left(K_z+\frac{EC_B}{2c}\right)d
       +K\left(\frac{C_B}{c}\right)^{1/p}d^{1/p}.
 \tag{S17}\label{eq:schur-matched-modulus}
\]
In particular, on a represented input block whose tangential inputs lie in
\(B_r^{n-1}\) and whose diameter is at most \(D>0\), a valid constant is
\[
 L_D=\left(K_z+\frac{EC_B}{2c}\right)D^{1-1/p}
           +K(C_B/c)^{1/p}.
\]
This is an alternative valid constant; it need not be smaller for every
choice of parameters.  For \(p>1\) and \(K_+,K_->0\), it reduces the
leading coefficient of \(d^{1/p}\) relative to the power allocation bound.
Condition \eqref{eq:schur-matched-output-jet} follows, for example,
by integrating
\(\|\partial_ty_+(z,t)-\partial_ty_-(z,t)\|\le pEt^{p-1}\).
Both reduced derivatives are the native Schur expressions in (S5),
evaluated at their corresponding tangentially corrected parameters.
\end{proposition}
\begin{proof}
Same-branch pairs already satisfy the stated bound.  For opposite
branches, suppose \(t\ge s\), compare first on the branch with parameter
\(t\), and then use \eqref{eq:schur-matched-output-jet} at \(s\):
\[
 \|y_+(z,t)-y_-(w,s)\|
 \le K_z\|z-w\|+K|t-s|+Es^p.
\]
If \(s\ge t\), interchange the branches and use the seam mismatch at
\(t\).  The opposite-branch estimate in the theorem gives
\(c(t^p+s^p)\le C_Bd\), whence
\[
 |t-s|\le(C_Bd/c)^{1/p},\qquad
 \min\{t,s\}^p\le C_Bd/(2c).
\]
Substitution proves the claim.
\end{proof}

\begin{proposition}[When the produced exponent is genuine]
\label{prop:schur-sharp-exponent}
Under the preceding corollary, fix a tangential coordinate \(z_0\).
Suppose that on at least one branch,
\[
 |n_\sigma(z_0,t)-\nu(z_0)|\le Ct^p,\qquad
 \|y_\sigma(z_0,t)-y_\sigma(z_0,0)\|\ge ct
\]
for \(0<t<t_0\), with \(C,c>0\).  Then no exponent greater than
\(1/p\) is valid on a graph neighborhood containing this arc and its
seam point.  In particular, for \(p>1\) the localized resolvent is not
calm there and hence is not locally Lipschitz.
\end{proposition}
\begin{proof}
Along this fixed tangential-input arc, an RL estimate of exponent
\(\theta\) would give \(ct\le LC^\theta t^{p\theta}\).
This is impossible as \(t\downarrow0\) when \(p\theta>1\).  Moreover,
\(2(u-u_0)=(\mathsf Xg-\mathsf Xg_0)+(\mathsf Yg-\mathsf Yg_0)\)
implies \(2\|u-u_0\|\ge ct-Ct^p\), which excludes calmness of the
resolvent for \(p>1\).
\end{proof}

\begin{remark}[Computable hypotheses and a finite atlas]
The Schur quantities in (S2)--(S6) are evaluated directly on graph
parametrizations.  No normal inverse or resolvent formula is a hypothesis.
If a graph branch is specified by an implicit active-manifold equation
\(G_\sigma(a,t,w)=0\), a certified invertible \(D_wG_\sigma\) produces
its derivatives through
\(D_{(a,t)}w=-(D_wG_\sigma)^{-1}D_{(a,t)}G_\sigma\).
Interval bounds on these matrices can therefore certify (S2)--(S6),
provided branch existence and the complete active-set list have been
established independently.

Several seam neighborhoods can be combined using a finite-chart chaining
argument, but that additional step still requires compatible
complete fibers on overlaps and a quantitative input-chain geometry.
The present theorem supplies a nontrivial \emph{local} derivative
certificate at a corank-one two-sided seam.  It is not a theorem for
arbitrary higher-codimension Whitney stratifications or arbitrarily many
branches meeting on the same side.
\end{remark}

\begin{proposition}[Native verification of the square-root operator]
\label{prop:schur-square-root-test}
Let \(\lambda=1\) and
\[
 F(\xi,y)=\{(-2\sqrt y,3y),(-2\sqrt y,-5y)\}\quad(y\ge0),
 \qquad F(\xi,y)=\varnothing\quad(y<0).
\]
Use the native, smooth graph parameters
\[
 g_+(a,t)=((a,t^2),(-2t,3t^2)),\qquad
 g_-(a,t)=((a,t^2),(-2t,-5t^2)),\quad t\ge0.
\]
On \(|a|\le R\), \(0\le t\le T\), the Schur data can be chosen as
\[
 \mu=1,\quad H=2,\quad B=0,\quad K_z=1,\quad
 h_+(t)=h_-(t)=4t^2,
\]
\[
 K_+=4\sqrt{1+T^2},\qquad K_-=4\sqrt{1+9T^2}.
\]
For every \(0<r<R-2T\), the complete graph over
\[
 U=(-r,r)\times(-4T^2,4T^2)
\]
is covered uniquely in Minty inputs.  On an input subblock of diameter
\(D\), the basic theorem yields
\[
 \gamma=\tfrac12,\qquad L_D=\sqrt D+\sqrt{8+40T^2}.
\]
The matched-jet version improves the asymptotic leading constant:
\[
 \gamma=\tfrac12,\qquad
 L_D^{\mathrm{jet}}=2\sqrt D+2\sqrt{1+9T^2}.
 \tag{S18}\label{eq:schur-square-root-L}
\]
In particular, along shrinking collars \(D,T\to0\),
\(L_D^{\mathrm{jet}}\to2\), which is the optimal asymptotic RL
constant.  These bounds are produced from graph derivatives without
assuming or solving a full inverse formula.
\end{proposition}
\begin{proof}
Direct calculation on the native graph gives
\[
 X_\pm(a,t)=(a-2t,\,\pm4t^2),\qquad
 Y_+(a,t)=(a+2t,-2t^2),\quad
 Y_-(a,t)=(a+2t,6t^2).
\]
Thus \(A_\pm=1\), \(b_\pm=-2\), and the signed normal Schur
derivatives are \(8t\).  The output Schur derivatives are
\((4,-4t)\) and \((4,12t)\), giving the asserted constants.
The common tangential correction is \(a=z+2t\).  Consequently the
two corrected output jets differ by \((0,-8t^2)\), so that
\eqref{eq:schur-matched-output-jet} holds with \(E=8\).  Formula
\eqref{eq:schur-square-root-L} follows from
Proposition~\ref{prop:schur-matched-jet}.

For completeness of the represented fibers, any full graph point with Minty
input in \(U\) has \(t=\sqrt y<T\) and
\(|a|=|z+2t|<r+2T<R\), and hence belongs to one of the displayed
graph parametrizations.  Conversely, the two monotone normal ranges
cover \((-4T^2,4T^2)\).  Finally, along \(z=0\), on the plus branch,
the input difference from the seam is \(4t^2\) and the output
difference is \(\sqrt{16t^2+4t^4}\).  Their RL ratio at exponent
\(1/2\) tends to \(2\), and every larger exponent fails.
\end{proof}

\begin{proposition}[Branchwise regularity does not replace orientation]
\label{prop:schur-collision-counterexample}
Even smooth, semialgebraic graph branches with a common seam, individually
one-to-one Minty maps, and individual optimal H\"older inverses need not
satisfy any all-pairs RL inequality after union.
\end{proposition}
\begin{proof}
In one input dimension, prescribe the two Minty-coordinate branches
\[
 X_+(t)=X_-(t)=t^2,\qquad
 Y_+(t)=t,\quad Y_-(t)=-t,\qquad t\ge0.
\]
They define operator graph branches through
\(u_\sigma=(X_\sigma+Y_\sigma)/2\),
\(v_\sigma=(X_\sigma-Y_\sigma)/(2\lambda)\).
Each branch is smooth and its Minty map is one-to-one, with
\(|Y_\sigma(t)-Y_\sigma(s)|\le|X_\sigma(t)-X_\sigma(s)|^{1/2}\).
Both branches agree at zero.  Their union is a smooth parabola in
\((X,Y)\)-coordinates; splitting it into its two open arcs and the
origin is a Whitney stratification.  Nevertheless, at any \(t>0\),
the two inputs agree, while the primal outputs differ by \(t\).
Thus every all-pairs RL inequality fails.  The missing hypothesis is
precisely opposite signed normal growth, not individual branch
smoothness or stratifiability.
\end{proof}

\begin{remark}[Uniformity in the tangential parameter is essential]
The smooth input map
\(X(a,t)=(a,t^3-at)\) has the cubic normal law \(t\mapsto t^3\)
at \(a=0\), but for each small \(a>0\) its normal coordinate is zero
at both \(t=0\) and \(t=\sqrt a\).  Taking
\(Y(a,t)=(a,t)\) again defines a legitimate operator graph with a
Minty collision.  A Taylor expansion at a single seam point cannot
replace the uniform Schur lower bound (S6).
\end{remark}

\paragraph{Relation to inverse-function machinery.}
The proof combines strong-monotonicity inversion in tangential coordinates,
a parameterized implicit-function calculation, and one-dimensional monotone
growth.  Its RLEB-specific output is the joint quantitative certificate for a
complete two-sided graph: a covered input collar, cross-branch collision
exclusion, and the constants in \eqref{eq:schur-explicit-L}.  Whitney
stratification alone is insufficient, as
Proposition~\ref{prop:schur-collision-counterexample} shows, and the residual
error bound must still be verified separately.
